\chapter{Painful Urination (Dysuria)}

\section*{Definition}
Pain, burning, or discomfort experienced during urination, typically indicating irritation or inflammation of the lower urinary tract (urethra, bladder, or prostate).

\section*{What Happens in Your Body}
Dysuria results from inflammation of the urethral or bladder mucosa, stimulating nociceptive nerve endings during urine passage. In women, the short urethra facilitates ascending bacterial infection (cystitis). In men, urethral irritation from infection (STI, prostatitis) or obstruction (BPH) is more common. Non-infectious causes include chemical irritation, interstitial cystitis, and atrophic vaginitis.

\section*{Causes}
\begin{itemize}
  \item Urinary tract infection (cystitis)
  \item Urethritis (gonorrhea, chlamydia, trichomonas)
  \item Prostatitis (acute or chronic bacterial)
  \item Interstitial cystitis (painful bladder syndrome)
  \item Atrophic vaginitis (postmenopausal)
  \item Chemical irritants (soaps, spermicides, hygiene products)
  \item Urethral stricture
  \item Bladder stones
  \item Medication side effects (cyclophosphamide)
  \item Sexually transmitted infections
\end{itemize}

\section*{When to See a Doctor}
See your doctor if:
\begin{itemize}
  \item Dysuria with fever, flank pain, blood in urine, vaginal/penile discharge, or inability to urinate
\end{itemize}

\section*{Related Symptoms}
\begin{itemize}
  \item Urinary frequency
  \item Urinary urgency
  \item Blood in urine
  \item Flank pain
  \item Fever
\end{itemize}


\section*{Self-Care / Home Management}
\begin{itemize}
  \item Drink fluids normally and enough to avoid dehydration; do not force large volumes, particularly if you have heart or kidney disease.
  \item Avoid potential irritants: caffeine, alcohol, spicy foods, acidic beverages, and nicotine, which can exacerbate bladder discomfort.
  \item Apply a warm compress to the suprapubic area to ease bladder discomfort; consider urinary alkalizers for symptom relief.
  \item Urinate frequently and completely; void immediately after intercourse to reduce bacterial colonisation risk.
\end{itemize}

\section*{How Your Doctor Will Evaluate You}
\begin{itemize}
  \item Urinalysis with microscopy for leukocytes, nitrites, blood, and bacteria; urine culture and sensitivity to identify the causative organism and antibiotic susceptibility.
  \item Physical examination including suprapubic palpation, costovertebral angle tenderness, and pelvic or prostate examination based on clinical context.
  \item Sexually transmitted infection screening (nucleic acid amplification tests for chlamydia, gonorrhoea, trichomonas) in sexually active individuals or when discharge is present.
  \item Imaging (renal ultrasound or CT urogram) if recurrent UTIs, haematuria, structural abnormalities, or stones are suspected.
\end{itemize}

\section*{Treatment Options}
\begin{itemize}
  \item Treatment depends on the cause and may include an antibiotic selected by a clinician, treatment for an STI, or management of prostatitis, atrophic tissue, stones, or another condition.
  \item Do not use leftover antibiotics or delay assessment when there is fever, flank pain, blood in the urine, discharge, pregnancy, or inability to urinate.
  \item Recurrent symptoms require evaluation rather than repeated self-treatment.
\end{itemize}

\section*{References}
\begin{thebibliography}{99}
\bibitem{painful_urination:dysuria1} Bremnor JD, Sadovsky R. ``Evaluation of dysuria in adults.'' Am Fam Physician. 2002;65(8):1589-1596.
\bibitem{painful_urination:dysuria2} Walsh CA, Moore KH. ``Dysuria.'' In: Evidence-Based Physical Diagnosis. 4th ed. Elsevier; 2018.
\bibitem{painful_urination:dysuria3} Gupta K, Hooton TM, Naber KG, et al. ``International clinical practice guidelines for the treatment of acute uncomplicated cystitis and pyelonephritis in women.'' Clin Infect Dis. 2011;52(5):e103-e120.
\bibitem{painful_urination:dysuria4} Michels TC, Sands JE. ``Dysuria: evaluation and differential diagnosis in adults.'' Am Fam Physician. 2015;92(9):778-786.
\bibitem{painful_urination:dysuria5} Hooton TM. ``Clinical practice. Uncomplicated urinary tract infection.'' N Engl J Med. 2012;366(11):1028-1037.
\bibitem{painful_urination:dysuria6} Schaeffer AJ, Nicolle LE. ``Urinary tract infections in adults.'' In: Campbell-Walsh Urology. 12th ed. Elsevier; 2021.
\end{thebibliography}
